% Part: conditional-logics
% Chapter: minimal-change-semantics
% Section: sphere-models

\documentclass[../../../include/open-logic-section]{subfiles}

\begin{document}

\olfileid{con}{min}{sph}

\olsection{ગોળા-નિદર્શો}

પ્રતિતથ્યાત્મક વિધાનો માટે ઔપચારિક અર્થવિચાર આપવાનો એક રસ્તો લૂઈસના
અનૌપચારિક વિશ્લેષણને ગાણિતિક સંરચનામાં ફેરવવાનો છે. ત્યારે
જગત~$w$ની આસપાસના ગોળા જગતોના ગણો બને છે. ગોળા એકબીજાની અંદર
હોવાથી~$w$ની આસપાસના જગતોના આ ગણો ઉપગણ-સંબંધ મુજબ રેખીય ક્રમમાં
હોવા જોઈએ.

\begin{defn}
  \emph{ગોળા-નિદર્શ} એ ત્રિક~$\mModel{M} = \tuple{W, O, V}$ છે,
  જેમાં $W$ જગતોનો અખાલી ગણ છે, $V\colon \PVar \to \Pow{W}$
  મૂલ્યાંકન છે, અને $O\colon W \to \Pow{\Pow{W}}$ દરેક જગત~$w$ને
  \emph{ગોળાઓનું તંત્ર}~$O_w$ આપે છે. દરેક $w$ માટે $O_w$ જગતોના
  ગણોનો ગણ છે અને તેણે નીચેની શરતો સંતોષવી જોઈએ:
  \begin{enumerate}
  \item $O_w$નું \emph{કેન્દ્ર}~$w$ છે: $\{w\} \in O_w$.
  \item $O_w$ \emph{અંતર્વિષ્ટ} છે: જ્યારે પણ $S_1$, $S_2 \in O_w$
    હોય, ત્યારે $S_1 \subseteq S_2$ અથવા $S_2 \subseteq S_1$;
    એટલે $O_w$નો ક્રમ~$\subseteq$ મુજબ રેખીય છે.
  \item $O_w$ તેના ગોળાઓના અખાલી સંગ્રહોના યોગગણ હેઠળ બંધ છે.
  \item $O_w$ તેના ગોળાઓના અખાલી સંગ્રહોના છેદગણ હેઠળ બંધ છે.
  \end{enumerate}
\end{defn}

$O_w$ પાછળનો સહજ વિચાર એવો છે કે~$w$ની ``આસપાસનાં'' જગતો
~$w$થી કેટલાં દૂર છે તે મુજબ સ્તરોમાં ગોઠવાય છે. સૌથી અંદરનો
ગોળો એકલું~$w$ છે, એટલે ગણ~$\{w\}$: બીજા કોઈ પણ ગોળામાંનાં
જગતો કરતાં~$w$ પોતે~$w$ની વધુ નજીક છે. જો $S \subsetneq S'$, તો
$S' \setminus S$માંનાં જગતો~$S$માંનાં જગતો કરતાં~$w$થી વધુ દૂર
છે: $S' \setminus S$ એ~$S$ અને~$S'$ની બહારનાં જગતો વચ્ચેનો
``સ્તર'' છે. ખાસ કરીને, ગોળાની બાહ્ય સપાટીની અંદરનાં બધાં જગતો
તેમાં આવે છે એવું સમજવું જોઈએ; ગોળા માત્ર અલગ અલગ સ્તરો નથી.

\begin{figure}
\begin{center}
\begin{tikzpicture}[layerwidth=1.5,scale=.6]\tiny
  \spheresystem{3}
  \propositionintersect{0}{3}{60}{6}
  \path (0,0) node[world] {$w$};
  \spherepos{-30}{2}{node[world] {$w_2$}}
  \spherepos{220}{2}{node[world] {$w_3$}}
  \spherepos{90}{2}{node[world] {$w_1$}}
  \spherepos[shift={(.1,0)}]{10}{3}{node[world] {$w_5$}}
  \spherepos[shift={(.1,0)}]{-10}{3}{node[world] {$w_6$}}
  \spherepos{225}{3}{node[world] {$w_4$}}
  \spherepos{20}{4}{node[world] {$w_7$}}
  \path (5,0) node[left] {$p$};
\end{tikzpicture}
\caption{ગોળા-નિદર્શનો આલેખ}
\ollabel{fig:sphere-model}
\end{center}
\end{figure}

\olref{fig:sphere-model}નો આલેખ એવા ગોળા-નિદર્શને અનુરૂપ છે જેમાં
$W = \{w, w_1, \dots, w_7\}$ અને $V(p) = \{w_5, w_6,
w_7\}$. સૌથી અંદરનો ગોળો $S_1 = \{w\}$ છે. $w$થી સૌથી નજીકનાં
જગતો $w_1, w_2, w_3$ છે, તેથી તેનાથી મોટો આગળનો ગોળો
$S_2 = \{w, w_1, w_2, w_3\}$ છે. વધુ દૂરનાં જગતો $w_4$, $w_5$,
$w_6$ છે, તેથી સૌથી બહારનો ગોળો
$S_3 = \{w, w_1, \dots, w_6\}$ છે. $w$ની આસપાસના ગોળાઓનું તંત્ર
$O_w = \{S_1, S_2, S_3\}$ છે. જગત~$w_7$,~$w$ની આસપાસના કોઈ
પણ ગોળામાં નથી. જ્યાં $p$ સત્ય છે એવાં સૌથી નજીકનાં જગતો $w_5$
અને~$w_6$ છે; તેથી $p$ને સ્વીકારતો સૌથી નાનો ગોળો~$S_3$ છે.

ગોળા-નિદર્શ~$\mModel M$માં જગત~$w$ પર સૂત્ર~$!A$ના સંતોષને,
$\mSat{M}{!A}[w]$, વ્યાખ્યાયિત કરવા મોડલ !!{formula}s માટેની
વ્યાખ્યામાં $!B \cif !C$ની શરત ઉમેરીએ છીએ:

\begin{defn}
  $\mSat{M}{!B \cif !C}[w]$ ત્યારે અને માત્ર ત્યારે જ્યારે નીચેની
  બેમાંથી કોઈ એક શરત સંતોષાય:
  \begin{enumerate}
  \item\ollabel{sphere-vac} દરેક $u \in \bigcup O_w$ માટે $\mSat/{M}{!B}[u]$, અથવા
  \item\ollabel{sphere-nonvac} કોઈ $S \in O_w$ માટે,
    \begin{enumerate}
      \item કોઈ $u \in
        S$ માટે $\mSat{M}{!B}[u]$, અને
      \item દરેક $v \in S$ માટે કાં તો
        $\mSat/{M}{!B}[v]$ અથવા $\mSat{M}{!C}[v]$.
    \end{enumerate}
  \end{enumerate}
\end{defn}

આ વ્યાખ્યા મુજબ $\mSat{M}{!B \cif !C}[w]$ ત્યારે અને માત્ર ત્યારે
જ્યારે પૂર્વપક્ષ~$!B$,~$w$ની આસપાસના બધા ગોળાઓમાં બધે અસત્ય હોય,
અથવા કોઈ ગોળા~$S$માં $!B$ સત્ય હોય અને એ ``$!B$-સ્વીકારતા''
ગોળાનાં બધાં જગતોમાં ભૌતિક શરતી વિધાન $!B \lif !C$ સત્ય હોય.
ધ્યાન આપો કે સહજ સમજૂતી પરથી અપેક્ષા રાખી શકાય તેમ વ્યાખ્યામાં
$S$ સૌથી અંદરનો $!B$-સ્વીકારતો ગોળો હોવો જોઈએ એવી શરત મૂકી
નથી. પરંતુ~\olref{sphere-nonvac}ની શરત કોઈ ગોળા~$S$ માટે સાચી
હોય, તો~$S$માં સમાયેલો અને પૂર્વપક્ષને સ્વીકારતો દરેક ગોળો પણ એ
શરત સંતોષે છે; તેથી સૌથી અંદરનો એવો ગોળો અસ્તિત્વમાં
હોય તો તે પણ સંતોષે છે.
\sourcecorrection{OLMD-109}{મૂળમાં Sમાં સમાયેલા દરેક ગોળા માટે શરત સાચી કહેવાઈ છે; નાનામાં B-સાક્ષી ન પણ હોય. અહીં ફક્ત B-સ્વીકારતા નાના ગોળા માટેનો સાચો દાવો રાખ્યો છે.}

એ પણ ધ્યાન આપો કે ગોળા-નિદર્શની વ્યાખ્યા મુજબ સૌથી અંદરનો
$!B$-સ્વીકારતો ગોળો \emph{હોવો જ જોઈએ} એવું નથી: $!B$ને સ્વીકારતા
ગોળાઓની અનંત શ્રેણી $S_1 \supsetneq S_2 \supsetneq \dots \supsetneq
\{w\}$ હોઈ શકે, અને તેથી સૌથી અંદરનો $!B$-સ્વીકારતો ગોળો ન હોય.
એવા કિસ્સામાં કોઈ~$i$ માટે $S_i$, $S_{i+1}$, \dots{} એ બધાં ગોળાઓમાં
$!B \lif !C$ સત્ય હોય ત્યારે અને માત્ર ત્યારે
$\mSat{M}{!B \cif !C}[w]$ છે.

\end{document}
